% !TeX program = pdflatex
\documentclass[12pt,a4paper]{article}
\usepackage[utf8]{inputenc}
\usepackage[T1]{fontenc}
\usepackage[polish]{babel}
\usepackage{lmodern}
\usepackage{geometry}
\usepackage{hyperref}
\usepackage{longtable}
\usepackage{enumitem}
\usepackage{csquotes}
\geometry{margin=2.5cm}
\hypersetup{colorlinks=true,linkcolor=blue,urlcolor=blue}

\title{Sprawozdanie – System RIMS (PostgreSQL)}
\author{}
\date{\today}

\begin{document}
\maketitle

\section{Jak powstał projekt (źródła i proces)}
Projekt zrealizowano w trzech krokach:
\begin{enumerate}
  \item \textbf{Modelowanie (Lucidchart)}: zaprojektowano schemat ERD i wyeksportowano go do pliku \texttt{MojeERD.csv}.
  \item \textbf{Analiza i specyfikacja (Gemini)}: plik CSV przeanalizowano w narzędziu \textbf{Gemini}; na tej podstawie powstały pliki \texttt{pomoc.md} (wskazówki dot. 3NF) oraz \texttt{polecenie.md} (wymagania i zakres).
  \item \textbf{Implementacja (Copilot + VS Code)}: w oparciu o powyższe materiały, z pomocą Copilota przygotowano schemat i dane w SQL, skrypty bash oraz dokumentację. Efektem pracy jest publiczne repozytorium i strona podglądu.
\end{enumerate}

Linki:
\begin{itemize}
  \item Repozytorium: \url{https://github.com/radekszewczyk01/RIMS_Database/}
  \item GitHub Pages (podgląd CSV/HTML): \url{https://radekszewczyk01.github.io/RIMS_Database/}
\end{itemize}

Pliki wejściowe i pochodzenie: \texttt{MojeERD.csv} (Lucidchart), \texttt{pomoc.md} i \texttt{polecenie.md} (Gemini), a także \texttt{pełna\_konwersacja.md} (zapis działań).

\section{Szczegóły techniczne}

\subsection{Struktura repozytorium}
Najważniejsze elementy: \texttt{MojeERD.csv} (eksport ERD z Lucidchart), \texttt{pomoc.md} (wytyczne 3NF), folder \texttt{sql/} (skrypty: m.in. \texttt{00\_all.sql}, \texttt{02\_schema.sql}, \texttt{03\_indexes\_views.sql}, \texttt{data\_and\_features.sql}), folder \texttt{scripts/} (\texttt{setup\_db.sh}, \texttt{show\_results.sh}) oraz \texttt{pełna\_konwersacja.md}.

\subsection{Model danych i normalizacja (3NF)}
Na bazie \texttt{pomoc.md} i ERD: dodano słowniki (\texttt{Dyscypliny}) i \texttt{Czasopisma}; znormalizowano łańcuch Wydawca$\rightarrow$Czasopismo$\rightarrow$Artykuł; PK to \texttt{SERIAL} (INT); kluczowe kolumny mają \texttt{NOT NULL}; relacje N:M są przez \texttt{Artykul\_Autor} i \texttt{Artykul\_ZrodloFinansowania}.

\subsection{Tabele (skrót)}
\begin{longtable}{p{0.25\linewidth}p{0.7\linewidth}}
\textbf{Wydawca} & id\_wydawcy PK, nazwa UNIQUE NOT NULL, kraj, srednia\_retrakcji, czy\_otwarty\_dostep \\
\textbf{Afiliacja} & id\_afiliacji PK, nazwa UNIQUE NOT NULL, kraj, miasto \\
\textbf{Dyscypliny} & id\_dyscypliny PK, nazwa UNIQUE NOT NULL \\
\textbf{Czasopisma} & id\_czasopisma PK, tytul UNIQUE NOT NULL, impact\_factor, id\_wydawcy FK, id\_dyscypliny FK \\
\textbf{Autor} & id\_autora PK, imie, nazwisko NOT NULL, orcid UNIQUE, id\_afiliacji FK \\
\textbf{ZrodloFinansowania} & id\_zrodla PK, nazwa NOT NULL, typ, kraj \\
\textbf{Artykul} & id\_artykulu PK, tytul NOT NULL, doi UNIQUE NOT NULL, rok\_publikacji NOT NULL, punkty\_mein, wspolczynnik\_rzetelnosci, data\_ostatniej\_aktualizacji, id\_czasopisma FK \\
\textbf{Artykul\_Autor} & id\_artykulu PK/FK, id\_autora PK/FK, kolejnosc\_autora \\
\textbf{Artykul\_ZrodloFinansowania} & id\_artykulu PK/FK, id\_zrodla PK/FK \\
\textbf{LogOceny} & id\_logu PK, id\_artykulu FK, data\_obliczenia NOT NULL, stary\_wspolczynnik, nowy\_wspolczynnik NOT NULL, id\_uzytkownika \\
\textbf{Cytowanie} & id\_cytowania PK, id\_cytujacego FK, id\_cytowanego FK, data\_zdarzenia \\
\textbf{ZarzutNierzetelnosci} & id\_zarzutu PK, typ\_nierzetelnosci NOT NULL, status, data\_zgloszenia, czy\_wycofany, id\_artykulu FK \\
\end{longtable}

\subsection{Relacje}
\begin{itemize}
  \item Wydawca 1—N Czasopisma —N Artykul
  \item Autor N—M Artykul (\texttt{Artykul\_Autor})
  \item ZrodloFinansowania N—M Artykul (\texttt{Artykul\_ZrodloFinansowania})
  \item Afiliacja 1—N Autor
  \item Artykul cytuje Artykul (\texttt{Cytowanie})
  \item Artykul 1—N ZarzutNierzetelnosci
  \item Artykul 1—N LogOceny
\end{itemize}

\subsection{Indeksy}
Indeksowano kolumny filtrujące i łączące: \texttt{Artykul(rok\_publikacji)}, \texttt{Artykul(data\_ostatniej\_aktualizacji)}, \texttt{Artykul(wspolczynnik\_rzetelnosci)}, \texttt{Afiliacja(kraj)} oraz pomocniczo \texttt{Artykul\_Autor(id\_autora)} i \texttt{(id\_artykulu)}. Wspomnianą w specyfikacji \texttt{data\_publikacji} zastąpiono praktycznymi polami.

\subsection{Widoki}
\begin{itemize}
  \item "\texttt{Widok\_Wycofane\_Artykuły\_Wydawcy}" – agregacja % wycofanych na wydawcę.
  \item \texttt{Widok\_Ryzykowne\_Finansowanie} – udział niskiego WR (< 0.2) i średni WR per źródło.
\end{itemize}

\subsection{Procedury}
\begin{itemize}
  \item \texttt{InsertNewArticle(tytul\_in, doi\_in, rok\_in, id\_czasopisma\_in, ...)} – dodaje wiersz do \texttt{Artykul}.
  \item \texttt{UpdateWR(article\_id\_in, new\_wr\_in, ...)} – aktualizuje WR oraz dopisuje wpis w \texttt{LogOceny}.
\end{itemize}

\subsection{Dane testowe i przykładowe wyniki}
Załadowano minimalnie 4 rekordy do każdej tabeli (więcej do \texttt{Artykul}).\newline
Przykłady:
\begin{itemize}
  \item \texttt{SELECT COUNT(*) FROM Artykul;} $\Rightarrow$ 24
  \item "Wycofane artykuły wydawcy" – SciPress (1/15), GlobalPub (1/3), AsiaAcad (1/3), EuroScience (0/3)
  \item "Ryzykowne finansowanie" – m.in. NCN, NSF, HorizonEU, PrivateTechFund, BioMedTrust
\end{itemize}

\subsection{Skrypty bash i uruchomienie}
	exttt{scripts/setup\_db.sh} uruchamia SQL przez stdin (\texttt{sudo -u postgres}); \texttt{scripts/show\_results.sh} wypisuje 4 zapytania (jeden widok w trybie \texttt{-XAtqc}).
Instrukcja skrótowa:
\begin{verbatim}
chmod +x scripts/*.sh
./scripts/setup_db.sh
./scripts/show_results.sh
\end{verbatim}

\subsection{Uwierzytelnianie i środowisko}
Lokalnie działa autoryzacja \emph{peer} dla roli \texttt{postgres} (stąd \texttt{sudo -u postgres}). Alternatywnie: \texttt{PGHOST}/\texttt{PGPORT}/\texttt{PGUSER}/\texttt{PGPASSWORD}.

\subsection{Gdzie są przechowywane dane}
W repozytorium dane do zasilenia znajdują się w pliku \texttt{sql/data\_and\_features.sql}.

Fizyczne pliki bazy leżą w katalogu danych PostgreSQL. Na tej maszynie:

\begin{itemize}
  \item \texttt{SHOW data\_directory;} $\Rightarrow$ \texttt{/var/lib/postgresql/16/main}
  \item przykładowa tabela \texttt{public.artykul}: \texttt{SELECT pg\_relation\_filepath('public.artykul');} $\Rightarrow$ \texttt{base/24576/24645}, czyli pełna ścieżka: \texttt{/var/lib/postgresql/16/main/base/24576/24645}
\end{itemize}
Tych plików nie modyfikuje się ręcznie — do pracy służą SQL/dumpy/restore.

% (sekcja celowo pominięta na życzenie: bez planów na przyszłość)

\end{document}